\chapter{Конкурентное сопровождение индекса}
\label{ch:maintenance}

\section{Сборка мусора в физическом порядке страниц}
% Логический обход дерева при вакуумировании даёт случайное чтение;
% физический порядок — последовательное.
% \commit{fe280694d0d}{Scan GiST indexes in physical order during VACUUM}.

\section{Удаление и переиспользование страниц}
% \commit{7df159a620b}{Delete empty pages during GiST VACUUM},
% \commit{9eb5607e699}{Refactor checks for deleted GiST pages}.

\section{Полные идентификаторы транзакций}
% Условие безопасного переиспользования страницы — «страница удалена достаточно
% давно». На 32-битных счётчиках это условие ломается при зацикливании.
% \commit{6655a7299d8}{Use full 64-bit XID for checking if a deleted GiST page is old enough}.
% Смежное: \commit{4ed8f0913bf}{Index SLRUs by 64-bit integers}.

\section{Снижение времени ожидания блокировок в обратном индексе}
% \commit{218f51584d5}{Reduce page locking in GIN vacuum} и последующие
% исправления взаимоблокировок \code{fd83c83d094}, \code{c6ade7a8cd3},
% \code{52ac6cd2d0c}. Опубликовано: JPCS 2018.
% Отдельно разобрать, почему первая версия породила три дефекта — это
% содержательный материал для главы о верификации.

\section{Предикатные блокировки}
% \commit{3ad55863e93}{Add predicate locking for GiST},
% \commit{0bef1c0678d}{Re-think predicate locking on GIN indexes}.
% Сериализуемая изоляция над обобщённым деревом: на каком уровне гранулярности
% брать предикатную блокировку, чтобы не порождать ложные конфликты.

\section{Упреждающее чтение при сопровождении}
% \commit{69273b818b1}{Use streaming read I/O in GiST vacuuming},
% \commit{e215166c9c8}{SP-GiST}, \commit{c5c239e26e3}{btree}.
% Связь с моделью §\ref{ch:models}.

\section{Экспериментальная оценка}

\section{Выводы по главе}
